%=============================================================================!
% 1.4  FROM ACCELERATION BACK TO MOTION                                       !
%=============================================================================!
\section{From acceleration back to motion}
\label{sec:mo-integrating}

Differentiation takes us from position to velocity and from velocity to
acceleration. To predict a motion we need to travel in the opposite
direction, from a known acceleration back to the position, and the tool
for this is integration.

\subsection*{Displacement as an area}

Divide the time from $0$ to $t$ into many short intervals of length
$\Delta t$, and let $t_k$ be the start of the $k$th interval. Over an
interval short enough that the velocity hardly changes, the object moves
through approximately $v(t_k)\,\Delta t$, a velocity multiplied by a time.
The total displacement is the sum of these small contributions,
\begin{equation*}
   x(t) - x(0) \approx \sum_k v(t_k)\,\Delta t ,
\end{equation*}
where $\sum_k$, a capital Greek sigma for `sum', means that the
expression after it is added up for every interval $k = 1, 2, 3, \dots$
in turn. Each term $v(t_k)\,\Delta t$ is the area of a thin rectangle of height
$v(t_k)$ and width $\Delta t$ standing on the time axis. As $\Delta t \to
0$ the approximation becomes exact and the sum becomes an integral, the
area under the graph of $v$. The same argument applied to the velocity,
whose rate of change is $a$, gives the companion result, and together
\begin{equation}
   v(t) = v(0) + \int_0^t a(t')\,\dd t', \qquad
   x(t) = x(0) + \int_0^t v(t')\,\dd t' .
   \label{eq:mo-integrals}
\end{equation}
These are the fundamental theorem of calculus, which states that
integrating a derivative from $0$ to $t$ returns the change in the
original function. The variable of integration is written $t'$ to
distinguish it from the upper limit $t$. Where $v$ is negative the area
counts as negative, so the integral measures the signed displacement. For
the ball, the area under $v(t)$ from the throw to the top of the flight is
therefore the greatest height that it reaches, shaded in
\cref{fig:mo-ball}(b).

\Cref{eq:mo-integrals} also shows what is needed to determine a motion.
Each integration introduces a starting value, the velocity $v(0)$ in the
first and the position $x(0)$ in the second, and these do not follow from
the acceleration. Two balls subject to the same acceleration but thrown
at different speeds therefore follow different paths. When the
acceleration is known at every instant, and the initial position and
velocity are given, the two integrals determine the motion uniquely.
Predicting a motion thus reduces to finding its acceleration, and
Chapter~2 shows that the acceleration is supplied by forces.

\subsection*{Constant acceleration}

When the acceleration is a constant $a$, both integrals in
\cref{eq:mo-integrals} can be evaluated directly. Writing $x_0 = x(0)$ and
$v_0 = v(0)$, the first gives the velocity,
\begin{equation*}
   v(t) = v_0 + \int_0^t a\,\dd t' = v_0 + at ,
\end{equation*}
and substituting this velocity into the second gives the position,
\begin{equation*}
   x(t) = x_0 + \int_0^t \bigl(v_0 + a t'\bigr)\,\dd t'
   = x_0 + \Bigl[\, v_0 t' + \tfrac12 a t'^2 \Bigr]_0^t
   = x_0 + v_0 t + \tfrac12 a t^2 .
\end{equation*}
Together,
\begin{equation}
   v(t) = v_0 + at, \qquad x(t) = x_0 + v_0 t + \tfrac12 a t^2 ,
   \label{eq:mo-constant}
\end{equation}
and \cref{eq:mo-ball} is the particular case $x_0 = 0$ and $a = -g$.

These two equations both contain the time, but many questions concern only
where an object goes, not when it gets there: how high does the ball
rise, for instance, or how fast is it moving at a given height? A relation
free of $t$ answers such questions directly. Provided that $a \neq 0$, the
first equation gives $t = (v - v_0)/a$, and substituting this into the
second eliminates the time:
\begin{align*}
   x - x_0 &= v_0\,\frac{v - v_0}{a} + \frac{a}{2}\,
      \frac{(v - v_0)^2}{a^2}
      && \text{substituting for $t$}\\
   &= \frac{2v_0(v - v_0) + (v - v_0)^2}{2a}
      && \text{cancelling $a$; common denominator $2a$}\\
   &= \frac{(v - v_0)\,(2v_0 + v - v_0)}{2a}
      && \text{taking out $(v - v_0)$}\\
   &= \frac{(v - v_0)(v + v_0)}{2a} = \frac{v^2 - v_0^2}{2a}
      && \text{$2v_0 - v_0 = v_0$; difference of squares.}
\end{align*}
Multiplying both sides by $2a$ and rearranging,
\begin{equation}
   v^2 = v_0^2 + 2a(x - x_0) .
   \label{eq:mo-v2}
\end{equation}

For the ball, $x_0 = 0$ and $a = -g$, so \cref{eq:mo-v2} reads $v^2 =
v_0^2 - 2gx$. The ball stops rising where $v = 0$, which gives $0 = v_0^2
- 2gx$ and hence the greatest height
\begin{equation*}
   x_{\max} = \frac{v_0^2}{2g} = \frac{12^2}{2\times9.80665}\,\si{\metre}
   = \SI{7.342}{\metre}.
\end{equation*}
The ball returns to the hand when $x = 0$ again. Setting $x = 0$ in
\cref{eq:mo-ball} and taking out the common factor $t$,
\begin{equation*}
   0 = v_0 t - \tfrac12 g t^2 = t\,\bigl(v_0 - \tfrac12 g t\bigr),
\end{equation*}
which is satisfied at the throw, $t = 0$, and again at $t = 2v_0/g =
\SI{2.447}{\second}$, twice the time to the top. At that moment
\cref{eq:mo-v2} with $x = 0$ gives $v^2 = v_0^2$, so $v = \pm v_0$; since
the ball is then falling, $v = -v_0$, and it returns at the speed with
which it was thrown.

When the acceleration depends on the position, as it will for an atom
pushed by its neighbours, the integrand of \cref{eq:mo-integrals} contains
the very motion we are trying to find. Except for special force laws,
we cannot evaluate those integrals in closed form. A computer then
advances the motion through many short intervals instead, and
\cref{sec:mo-taylor} develops the mathematics that makes this possible,
once the trigonometric functions it relies on are in place.

\notebook{01}{the area under $v(t)$ up to the cursor}
